This technical note documents the agricultural contribution to PakEmis using a common livestock, crop, fertiliser and nitrogen-flow account for greenhouse gases and air pollutants. The supported national core covers July 2015--June 2025. FY2023--24 remains the detailed reference year; earlier long calendar-year diagnostics are separate products and do not define the common fiscal period.

The reference TAN-coherent configuration produces \SI{\SectorRefGHG}{\mega\tonne} \COe{} on AR5 GWP-100, \SI{\SectorRefNHthree}{\kilo\tonne} \NHthree{}, \SI{\SectorRefNOx}{\kilo\tonne} NO$_2$-equivalent \NOx{} and \SI{\SectorRefPMfine}{\kilo\tonne} primary PM$_{2.5}$ in FY2023--24. FY2024--25 uses observed fiscal crop and nutrient controls and the consistent census-anchored livestock reconstruction. Native gas masses, pollutant masses and GWP aggregates remain separate.

Species- and stage-resolved TAN losses drive both \NHthree{} and the surviving manure pools used for \NtwoO{}. An IPCC volatilisation configuration is retained separately; it is not mixed probabilistically with the TAN reference. Conditional uncertainty replays shared herd, crop, nitrogen and applicable fate-factor draws, while management shares, factor applicability and missing sources remain distinct limitations. Direct comparison with the publicly released HTAPv3.2 monthly country tables matches exact July--June periods and records shared REAS/EDGAR inputs. The release provides national activity detail and process-matched geographic support beyond a simple clipped global inventory.
